\documentclass[a4paper,fontset=none]{ctexart}

\usepackage{anyfontsize}
\usepackage{amsmath,amssymb,mathrsfs,wasysym}
\usepackage{autobreak}
\allowdisplaybreaks
\usepackage[T1]{fontenc}
\usepackage{textcomp}
\usepackage{amsthm}
\usepackage[libertine,vvarbb]{newtxmath}
\usepackage[scr=rsfso,frak=euler,bb=ams]{mathalfa}
\usepackage{bm}
\usepackage{indentfirst}
\setlength{\parindent}{0em}
\usepackage{parskip}
\usepackage{geometry}
\geometry{top=3cm,bottom=3cm,left=2cm,right=2cm}
\usepackage{fontspec}
\setmainfont{Linux Libertine O}
\setsansfont{Linux Biolinum O}
\setmonofont{JetBrains Mono}
\setCJKmainfont{SimSun}[BoldFont=Noto Sans CJK SC Medium]

\begin{document}

    \title{\textbf{电动力学作业}}
    \author{1900011604\ \ 张植竣\ \ 北京大学}
    \date{2021年4月20日}
    \maketitle

    \section*{第四章}

    \begin{itemize}

    \item[2.]
    已知$n=1$，$n'=\sqrt{2}$，入射角$i=45^\circ$，由Snell定律$n\sin i=n'\sin r$得折射角$r=30^\circ$。
    
    由Fresnel公式：
    \[
        \frac{E'_\perp}{E_\perp} = \frac{2n\cos i}{n\cos i + n'\cos r} = \sqrt{3}-1
    \]
    \[
        \frac{E''_\perp}{E_\perp} = \frac{n\cos i-n'\cos r}{n\cos i + n'\cos r} = \frac{\sqrt{2}-\sqrt{6}}{\sqrt{2}+\sqrt{6}}
    \]
    则有：
    \[
        R = \left(\frac{E''_\perp}{E_\perp}\right)^2 = \frac{2-\sqrt{3}}{2+\sqrt{3}}
    \]
    \[
        T = 1 - R = \frac{n'\cos r}{n\cos i}\left(\frac{E'_\perp}{E_\perp}\right)^2 = \frac{2\sqrt{3}}{2+\sqrt{3}}
    \]

    \item[3.]
    由Snell定律$n\sin i=n'\sin r$得$\sin r=1.152>1$，$r$无实数解，说明发生了全反射，此时$r$取复数。

    沿表面传播的折射波为：
    \[
        \mathrm{e}^{i\bm{k}'\cdot\bm{x}} = \mathrm{e}^{-k'z|\cos r|} \cdot \mathrm{e}^{ik'x\tfrac{\sin i}{\sin i_0}}
    \]
    其中$\sin i_0=\dfrac{n'}{n}$

    设波速为$v'_p$，则：
    \[
        \frac{\omega}{v'_p} = k'\left(\frac{\sin i}{\sin i_0}\right) = k'\sin i\left(\frac{n}{n'}\right) = k\sin i
    \]
    \[
        v_p = \frac{v_p}{\sin i} = \frac{c}{n\sin i}
    \]
    使用$n=1.33\approx\dfrac{1}{\sin^2 i}$的近似，可以得到：
    \[
        v_p \approx \dfrac{c}{\sin i}\sin^2 i = \frac{\sqrt{3}}{2}c
    \]
    穿透深度：
    \[
        \kappa^{-1} = \frac{1}{k'|\cos r|} = k'^{-1}\left[\left(\frac{\sin i}{\sin i_0}\right)^2-1\right]^{-\tfrac{1}{2}} = \frac{\lambda'}{2\pi}\left[\left(\frac{\sin i}{\sin i_0}\right)^2-1\right]^{-\tfrac{1}{2}}
    \]
    代入$\lambda'=6.28\times10^{-5}\,\mathrm{cm}$得：$\kappa^{-1}=1.75\times10^{-5}\,\mathrm{cm}$

    \item[4.]
    介质内的Maxwell方程为：
    \[
        \nabla\cdot\bm{D} = 0,\quad \nabla\times\bm{E} + \frac{\partial\bm{B}}{\partial\bm{t}} = 0
    \]
    \[
        \nabla\cdot\bm{B} = 0,\quad \nabla\times\bm{H} - \frac{\partial\bm{D}}{\partial\bm{t}} = 0
    \]
    由题意：$\bm{D}\neq\epsilon\bm{E}$，$\bm{B}=\mu\bm{H}$。

    设：
    \[
        \bm{D} = \bm{D}_0\mathrm{e}^{i(\bm{k}\cdot\bm{x}-\omega t)},\quad \bm{E} = \bm{E}_0\mathrm{e}^{i(\bm{k}\cdot\bm{x}-\omega t)}
    \]
    \[
        \bm{B} = \bm{B}_0\mathrm{e}^{i(\bm{k}\cdot\bm{x}-\omega t)},\quad \bm{H} = \bm{H}_0\mathrm{e}^{i(\bm{k}\cdot\bm{x}-\omega t)}
    \]

    \item[7.]
    透入深度：
    \[
        \delta = \sqrt{\frac{2}{\mu\omega\sigma}} = \sqrt{\frac{2}{\mu_0 \mu_r \cdot (2\pi\nu) \cdot \sigma}} = (\mu_0 \mu_r \pi\nu\sigma)^{-\tfrac{1}{2}}
    \]
    $\nu=50\,\mathrm{Hz}$时，$\delta=(\mu_0 \mu_r \pi\nu\sigma)^{-\tfrac{1}{2}}=71.176\,\mathrm{m}$；

    $\nu=10^6\,\mathrm{Hz}$时，$\delta=(\mu_0 \mu_r \pi\nu\sigma)^{-\tfrac{1}{2}}=0.5033\,\mathrm{m}$；

    $\nu=10^9\,\mathrm{Hz}$时，$\delta=(\mu_0 \mu_r \pi\nu\sigma)^{-\tfrac{1}{2}}=0.0159\,\mathrm{m}$。

    \item[8.]
    设入射面为$xz$平面，界面为$z=0$，$z>0$为真空部分，$z<0$为介质部分。

    入射波：$\bm{k}_1=k_1\sin\theta_1\hat{x}+k_1\cos\theta_1\hat{z}$。

    透射波为复数$\bm{k}_2=\bm{a}+i\bm{b}$，其中$k_{2x}=a_x+ib_x$，$k_{2z}=a_z+ib_z$。

    由于$z=0$平面上入射波与透射波相因子相等：
    \[
        k_1\sin\theta_1 = a_x + ib_x
    \]
    则有：
    \[
        a_x = k_1\sin\theta_1,\quad b_x = 0
    \]
    波矢大小的关系为：
    \[
        k_1 = \frac{\omega n}{c} = \frac{\omega}{c}
    \]
    \[
        k_2 = \frac{\omega n'}{c} = \frac{\omega}{c}\sqrt{(\epsilon+i\frac{\sigma}{\omega})\mu c^2} = \omega\sqrt{\mu(\epsilon+i\frac{\sigma}{\omega})}
    \]
    注意到由于$\bm{k}_2=a_x\hat{x}+(a_z+ib_z)\hat{z}$，有：
    \[
        k_2^2 = a_x^2 + (a_z^2 + 2ia_z b_z - b_z^2) = (a_x^2 + a_z^2 - b_z^2) + 2ia_z b_z
    \]
    与$k_2^2=\omega^2\mu\epsilon+i\omega\mu\sigma$联立得：
    \[
        a_x^2 + a_z^2 - b_z^2 = \omega^2\mu\epsilon,\quad 2a_z b_z = \omega\mu\sigma
    \]
    代入$a_x^2=(k_1\sin\theta_1)^2=\dfrac{\omega^2}{c^2}\sin^2\theta_1$，解得：
    \[
        a_z^2 = \frac{1}{2}(\omega^2\mu\epsilon-\frac{\omega^2}{c^2}\sin^2\theta_1) + \frac{1}{2}\sqrt{(\omega^2\mu\epsilon-\frac{\omega^2}{c^2}\sin^2\theta_1)^2+(\omega\mu\sigma)^2}
    \]
    \[
        b_z^2 = -\frac{1}{2}(\omega^2\mu\epsilon-\frac{\omega^2}{c^2}\sin^2\theta_1) + \frac{1}{2}\sqrt{(\omega^2\mu\epsilon-\frac{\omega^2}{c^2}\sin^2\theta_1)^2+(\omega\mu\sigma)^2}
    \]
    介质中的透射波有如下形式：
    \[
        \mathrm{e}^{-k_z z+ik_x x}=\mathrm{e}^{-(a_z+ib_z)z+ia_x x}
    \]
    相速度：$v_p=\dfrac{\omega}{a}=\dfrac{\omega}{\sqrt{a_x^2+a_z^2}}$，透入深度：$\delta=\dfrac{1}{b_z}$

    若导电介质为金属：$\dfrac{\sigma}{\omega\epsilon}\ll1$

    注意到$a_x^2=\dfrac{\omega^2}{c^2}\sin^2\theta_1=\omega^2\mu_0\epsilon_0\sin^2\theta_1\sim\omega^2\mu\epsilon$，则$\omega\mu\sigma\gg\omega\mu\epsilon-\dfrac{\omega^2}{c^2}\sin^2\theta_1$。
    此时有：
    \[
        a_x \ll a_z,\quad a_z \approx b_z \approx \sqrt{\frac{\omega\mu\sigma}{2}}
    \]
    则相速度和透入深度为：
    \[
        v_p = \frac{\omega}{a_z} = \sqrt{\frac{2\omega}{\mu\sigma}},\quad \delta = \frac{1}{b_z} = \sqrt{\frac{2}{\omega\mu\sigma}}
    \]

    \item[9.]
    波导管一端封闭，不具有空间对称性，则只考虑时间因子$\bm{E}=\bm{E_0}(x,y,z)\mathrm{e}^{-i\omega t}$。

    $\bm{E}_0$满足：$\nabla^2\bm{E}_0+k^2\bm{E}_0=0$，边界条件：$\bm{n}\times\bm{E}_0=0$，Maxwell方程：$\nabla\cdot\bm{E}_0=0$。

    $\bm{E}_0$的三个分量$E_x$、$E_y$、$E_z$也分别满足：
    \[
        \nabla^2E_i + k^2E_i = \nabla^2E_i + (k_x^2+k_y^2+k_z^2)E_i = 0,\quad i = x,y,z
    \]
    由分离变量法：
    \[
        E_i = \sum_{j} (C_j\cos k_j x_j + D_j\sin k_j x_j),\quad i,j = x,y,z
    \]
    设波导管在$x$方向宽度为$a$，在$y$方向宽度为$b$。

    由$\nabla\cdot\bm{E}_0=0$及边界条件$\bm{n}\times\bm{E}_0=0$得：
    \[
        \frac{\partial E_x}{\partial x} = 0,\quad E_y = E_z = 0,\quad (x=0,x=a)
    \]
    \[
        \frac{\partial E_y}{\partial y} = 0,\quad E_x = E_z = 0,\quad (y=0,y=b)
    \]
    \[
        \frac{\partial E_z}{\partial z} = 0,\quad E_x = E_y = 0,\quad (z=0)
    \]
    可以确定：
    \[
        \left\{
        \begin{aligned}
            & E_x = A_x\cos\frac{m\pi x}{a}\sin\frac{n\pi y}{b}\sin(k_z z) \\
            & E_y = A_y\sin\frac{m\pi x}{a}\cos\frac{n\pi y}{b}\sin(k_z z)\quad m,n\in\{0\}\cup\mathbb{N}^+ \\
            & E_z = A_z\sin\frac{m\pi x}{a}\sin\frac{n\pi y}{b}\cos(k_z z) \\
        \end{aligned}
        \right.
    \]
    而$k_z$只能通过$k_x^2+k_y^2+k_z^2=k^2=\left(\dfrac{\omega}{c}\right)^2$确定：
    \[
        k_z=\sqrt{\left(\frac{\omega}{c}\right)^2-\left(\frac{m\pi}{a}\right)^2-\left(\frac{n\pi}{b}\right)^2}
    \]
    又因为$\nabla\cdot\bm{E}=\dfrac{\partial E_x}{\partial x}+\dfrac{\partial E_y}{\partial y}+\dfrac{\partial E_z}{\partial z}=0$，代入得：$A_x k_x+A_y k_y + A_z k_z=0$，即：
    \[
        A_x\frac{m\pi}{a} + A_y\frac{m\pi}{b} + A_z\sqrt{\left(\dfrac{\omega}{c}\right)^2-\left(\dfrac{m\pi}{a}\right)^2-\left(\dfrac{n\pi}{b}\right)^2} = 0
    \]
    可见$A_x$、$A_y$、$A_z$中只有两个是独立的，即对每一组$(m,n)$，有两种独立的波模。

    \item[14.]
    由于$x$方向均匀，此时只考虑$y$方向的电场分布：$\bm{E}=\bm{E}_0(y)\mathrm{e}^{i(k_z z-\omega t)}$

    $\bm{E}_0$满足：$\nabla^2\bm{E}_0+k^2\bm{E}_0=0$，边界条件：$\bm{n}\times\bm{E}_0=0$，Maxwell方程：$\nabla\cdot\bm{E}_0=0$。

    $\bm{E}_0$的三个分量$E_x$、$E_y$、$E_z$也分别满足：
    \[
        \nabla^2E_i + k^2E_i = \frac{\mathrm{d}^2E_i}{\mathrm{d}y^2} + k_y^2E_i = 0,\quad i = x,y,z
    \]
    由分离变量法：
    \[
        E_i = C_i\cos k_y y + D_j\sin k_y y,\quad i = x,y,z
    \]
    由$\nabla\cdot\bm{E}_0=0$及边界条件$\bm{n}\times\bm{E}_0=0$得：
    \[
        \frac{\partial E_y}{\partial y} = 0,\quad E_x = E_z = 0,\quad (y=0,y=b)
    \]
    可以确定：
    \[
        \left\{
        \begin{aligned}
            & E_x = A_x\sin\frac{m\pi y}{b} \\
            & E_y = A_y\cos\frac{m\pi y}{b}\quad m\in\{0\}\cup\mathbb{N}^+ \\
            & E_z = A_z\sin\frac{m\pi y}{b} \\
        \end{aligned}
        \right.
    \]
    而$k_z$只能通过$k_y^2+k_z^2=k^2=\left(\dfrac{\omega}{c}\right)^2$确定：
    \[
        k_z=\sqrt{\left(\frac{\omega}{c}\right)^2-\left(\frac{m\pi}{b}\right)^2}
    \]
    又因为$\nabla\cdot\bm{E}=\dfrac{\partial E_x}{\partial x}+\dfrac{\partial E_y}{\partial y}+\dfrac{\partial E_z}{\partial z}=0$，代入得：$A_y k_y - iA_z k_z=0$，即：
    \[
        A_y\frac{m\pi}{b} - iA_z\sqrt{\left(\frac{\omega}{c}\right)^2-\left(\frac{m\pi}{b}\right)^2} = 0
    \]
    而$\dfrac{\partial E_x}{\partial x}=0$，$A_x$是独立的。
    
    可见$A_x$、$A_y$、$A_z$中只有两个是独立的，即对每一个$m$值，有两种独立的波模。

    截止频率为$\omega_c=k_y c=\dfrac{m\pi}{b}c$

    \item[15.]
    对于边长为$a$、$b$、$c$的谐振腔，由分离变量法可得：
    \[
        \left\{
        \begin{aligned}
            & E_x = A_x\mathrm{e}^{-i\omega t}\cos\frac{m\pi x}{a}\sin\frac{n\pi y}{b}\sin\frac{p\pi z}{c} \\
            & E_y = A_y\mathrm{e}^{-i\omega t}\sin\frac{m\pi x}{a}\cos\frac{n\pi y}{b}\sin\frac{p\pi z}{c}\quad m,n,p\in\{0\}\cup\mathbb{N}^+ \\
            & E_z = A_z\mathrm{e}^{-i\omega t}\sin\frac{m\pi x}{a}\sin\frac{n\pi y}{b}\cos\frac{p\pi z}{c} \\
        \end{aligned}
        \right.
    \]
    且$k_x=\dfrac{m\pi}{a}$、$k_y=\dfrac{n\pi}{b}$、$k_z=\dfrac{p\pi}{c}$满足：
    \[
        k_x^2 + k_y^2 + k_z^2 = \left(\frac{\omega}{c}\right)^2
    \]
    \[
        A_x k_x+A_y k_y + A_z k_z = 0
    \]
    则电场能量密度的时间平均值为：
    \[
        \frac{1}{2}\epsilon_0|\bm{E}|^2 = \frac{1}{4}\epsilon_0(A_x^2 + A_y^2 + A_z^2)
    \]
    将电场的$x$、$y$、$z$分量代入$\bm{B}=\dfrac{1}{i\omega}\nabla\times\bm{E}$得：
    \[
        \left\{
        \begin{aligned}
            & B_x = \frac{1}{i\omega}(k_y A_z - k_z A_y)\mathrm{e}^{-i\omega t}\cos\frac{m\pi x}{a}\sin\frac{n\pi y}{b}\sin\frac{p\pi z}{c} \\
            & B_y = \frac{1}{i\omega}(k_z A_x - k_x A_z)\mathrm{e}^{-i\omega t}\sin\frac{m\pi x}{a}\cos\frac{n\pi y}{b}\sin\frac{p\pi z}{c}\quad m,n,p\in\{0\}\cup\mathbb{N}^+ \\
            & B_z = \frac{1}{i\omega}(k_y A_x - k_x A_y)\mathrm{e}^{-i\omega t}\sin\frac{m\pi x}{a}\sin\frac{n\pi y}{b}\cos\frac{p\pi z}{c} \\
        \end{aligned}
        \right.
    \]
    磁场密度的时间平均值为：
    \begin{align*}
        \frac{1}{2\mu_0}|\bm{B}|^2
        &= \frac{1}{4\mu_0\omega^2}(k_y^2A_z^2+A_y^2k_z^2-2k_y A_z A_y k_z+k_z^2A_x^2+A_z^2k_x^2-2k_z A_x A_z k_x+k_y^2A_x^2+A_y^2k_x^2-2k_y A_x A_y k_x) \\
        &= \frac{1}{4\mu_0\omega^2}[(k_x^2+k_y^2+k_z^2)(A_x^2+A_y^2+A_z^2) + (A_x k_x+A_y k_y + A_z k_z)^2] \\
        &= \frac{1}{4\mu_0\omega^2} \cdot \frac{\omega^2}{c^2}(A_x^2+A_y^2+A_z^2) \\
        &= \frac{1}{4}\epsilon_0(A_x^2+A_y^2+A_z^2)
    \end{align*}
    故电场能量和磁场能量对时间的平均值总相等。

    \end{itemize}

    \subsection*{补充题}

    \begin{itemize}

        \item[1.]
        设振幅为$E_0$，则$\bm{E}$可以由$\hat{x}$和$\hat{y}$方向的线偏振波叠加，解析式为：
        \[
            \bm{E} = \frac{1}{\sqrt{2}}E_0\mathrm{e}^{i(kz-\omega t)}(\hat{x}+\hat{y})
        \]
        将其分解为圆偏振波的基：
        \[
            \hat{e}_R = \frac{1}{\sqrt{2}}(\hat{x} - i\hat{y})\quad \hat{e}_L = \frac{1}{\sqrt{2}}(\hat{x} + i\hat{y})
        \]
        可得：
        \begin{align*}
            \bm{E}
            &= \frac{1}{\sqrt{2}}E_0\mathrm{e}^{i(kz-\omega t)}(\hat{x}+\hat{y}) \\
            &= \frac{1}{2}E_0\mathrm{e}^{i(kz-\omega t)}[(1+i)\hat{e}_R + (1-i)\hat{e}_L] \\
            &= \frac{1+i}{2}E_0\mathrm{e}^{i(kz-\omega t)}\hat{e}_R + \frac{1-i}{2}E_0\mathrm{e}^{i(kz-\omega t)}\hat{e}_L
        \end{align*}

        \item[2.(1)]
        设入射波的电场强度振幅为$\bm{E}$，折射波为$\bm{E}'$，反射波为$\bm{E}''$，且入射介质为$(\epsilon_1,\mu_0)$，折射介质为$(\epsilon_2,\mu_0)$，则能流密度为：
        \[
            \bm{S} = \sqrt{\frac{\epsilon_1}{\mu_0}}|\bm{E}|^2\hat{k}
        \]
        \[
            \bm{S}' = \sqrt{\frac{\epsilon_2}{\mu_0}}|\bm{E'}|^2\hat{k}'
        \]
        \[
            \bm{S}'' = \sqrt{\frac{\epsilon_1}{\mu_0}}|\bm{E''}|^2\hat{k}''
        \]

        \item[2.(2)]
        能量守恒要求：
        \[
            (\bm{S} + \bm{S}'' - \bm{S}') \cdot \bm{n} = 0
        \]

        \item[2.(3)]
        全反射时，没有折射波，能量守恒要求：
        \[
            (\bm{S} + \bm{S}'') \cdot \bm{n} = 0,\quad S' \cdot \bm{n} = 0
        \]

    \end{itemize}

\end{document}
